\begin{ExerciseBlock}
\ExerciseHeading{04}{Ballistisches Pendel}{two_particle_systems_ballistic_pendulum}
In der Abbildung ist ein sogenanntes ballistisches Pendel dargestellt. Eine Masse $M = 2,5\,$kg ist an einem Seil der Länge $l = 1,0\,$m aufgehängt. Eine Masse $m = 30\,$g fliegt mit einer Geschwindigkeit $v_0$ horizontal auf die Masse $M$ zu, und bleibt in ihr stecken. Das Pendel wird um $\varphi = 20^\circ$ ausgelenkt. Im Folgenden sollen die Massen als Punktmassen behandelt werden. Es soll die Geschwindigkeit $v_0$ berechnet werden. Dazu arbeiten wir rückwärts:

\begin{center}
	\CachedTikzFigure{content/resources/chapters/two_particle_systems/tikz/ballistic_pendulum.tex}
\end{center}

\begin{enumerate}
	\item Leiten Sie eine Formel für die Höhe $h$ her, um die das Pendel ausgelenkt wird. Die Formel sollte nur von $l$ und $\varphi$ abhängen. (Dies ist reine Trigonometrie.)

	\item Leiten Sie nun eine Gleichung für die Geschwindigkeit $v_1$ her, um diese Auslenkung zu
erzielen.

	\item Leiten Sie schließlich eine Formel für die Geschwindigkeit $v_0$ her und berechnen Sie diese.

\end{enumerate}
\end{ExerciseBlock}
\UseExerciseSolution{two_particle_systems_ballistic_pendulum}
